\documentclass[11pt,a4paper]{article}
\usepackage[T1]{fontenc}
\usepackage[utf8]{inputenc}
\usepackage[czech]{babel}
\usepackage{lmodern,microtype,amsmath,amssymb,geometry,enumitem,tcolorbox}
\geometry{top=18mm,bottom=18mm,left=22mm,right=22mm}
\pagestyle{empty}
\setlength{\parindent}{0pt}
\setlength{\parskip}{4pt}
\newtcolorbox{stanoviste}[1]{colback=white,colframe=black,boxrule=0.45pt,arc=0.8mm,colbacktitle=white,coltitle=black,title={#1},fonttitle=\bfseries}
\begin{document}
\begin{center}{\LARGE\bfseries Matematická únikovka}\par\medskip{\large 136. hodina -- finále prvního ročníku}\end{center}
\textbf{Skupina:} \rule{75mm}{0.4pt}\hfill\textbf{Třída:} \rule{24mm}{0.4pt}

\textbf{Pravidla:} Každé ze šesti stanovišť poskytuje \emph{jednu číslici}. Číslice zapisujte podle pořadí stanovišť. Před odevzdáním vždy ověřte výsledek. U matematických řešení uvádějte podmínky.

\begin{stanoviste}{Stanoviště 1 -- Logická karta}
V množině $\{1;2;\dots;9\}$ spočítejte čísla, která jsou \emph{sudá nebo dělitelná třemi, ale ne obojí zároveň}. Číslicí je počet těchto čísel.\par\vspace{8mm}\end{stanoviste}
\begin{stanoviste}{Stanoviště 2 -- Ztracený jmenovatel}
Najděte $x$ z rovnice $\dfrac{x^2-9}{x-3}=7$. Pozor na podmínku. Číslicí je získaný kořen.\par\vspace{8mm}\end{stanoviste}
\begin{stanoviste}{Stanoviště 3 -- Přímka bez předpisu}
Přímka prochází body $[1;8]$ a $[3;4]$. Určete předpis funkce a její hodnotu $f(2)$. Číslicí je tato hodnota.\par\vspace{8mm}\end{stanoviste}
\begin{stanoviste}{Stanoviště 4 -- Dvě neznámé}
Vyřešte $\begin{cases}x+y=9,\\2x-y=3.\end{cases}$ Číslicí je hodnota $y$.\par\vspace{8mm}\end{stanoviste}
\begin{stanoviste}{Stanoviště 5 -- Modul jako vzdálenost}
Rovnice $|x-4|=3$ má dva kořeny. Určete je. Číslicí je jejich \emph{součet}.\par\vspace{8mm}\end{stanoviste}
\begin{stanoviste}{Stanoviště 6 -- Poslední parabola}
Vyřešte $x^2-5x+6=0$. Číslicí je \emph{součin} reálných kořenů.\par\vspace{8mm}\end{stanoviste}
\bigskip
\begin{center}
{\Large\bfseries Finální kód}\par\vspace{3mm}
\begin{tabular}{|p{1.05cm}|p{1.05cm}|p{1.05cm}|p{1.05cm}|p{1.05cm}|p{1.05cm}|}\hline\rule{0pt}{12mm} & & & & & \\\hline\end{tabular}
\end{center}
\end{document}
